\chapter{Gaussian variation and the population linear projection}
\label{gaussian:chapter}

A fitness relationship specifies an expected output at each breeding value.
A variance calculation additionally needs to know how often those breeding
values occur. This chapter derives the calculation for a Gaussian population,
then shows exactly which steps change when the population is not Gaussian.
The regression interpretation will follow from minimising a population mean
squared error. It is not an assumption that the local tangent is the fitted
line.

\section{The function and the population are separate inputs}
\label{gaussian:quadratic_target}

Let the breeding-value vector be $\vect A\in\R^n$, with population mean
$\vect\mu$ and covariance $\mat G$. Write
\begin{equation}
\vect X=\vect A-\vect\mu,
\qquad \E[\vect X]=\vect0,
\qquad \E[\vect X\vect X^{\mathsf T}]=\mat G.
\end{equation}
Breeding values use the population-specific average effects appropriate to
the trait and reference population. The following mathematics concerns the
distribution of those values; it does not derive the effects from a genetic
architecture.

Suppose log conditional expected reproductive output, $\ell_A$, is twice
continuously differentiable near $\vect\mu$. Its quadratic expansion is
\begin{align}
\ell_A(\vect\mu+\vect x)
 &=q(\vect x)+R(\vect x),\nonumber\\
q(\vect x)
 &=\ell_0+\vect b^{\mathsf T}\vect x
       +\frac12\vect x^{\mathsf T}\mat H\vect x,
\label{gaussian:quadratic_definition}
\end{align}
where $\ell_0=\ell_A(\vect\mu)$, $\vect b=\nabla\ell_A(\vect\mu)$ and
$\mat H=\nabla^2\ell_A(\vect\mu)$. Equality of the mixed second derivatives
makes $\mat H$ symmetric. The factor $1/2$ is the Taylor coefficient. Thus,
in two traits, the quadratic term is
\[
\frac12H_{11}x_1^2+H_{12}x_1x_2+\frac12H_{22}x_2^2.
\]
The off-diagonal contribution appears twice in the matrix product before
the factor $1/2$ is applied.

The function $q$ is the same quadratic approximation denoted
$\widetilde\ell_A$ elsewhere in the companion. Under the differentiability
assumption, $R(\vect x)=o(\norm{\vect x}^2)$ as $\vect x\to\vect0$.
This is a local statement. A Gaussian population has unbounded support, so
a local remainder statement alone does not establish that $R(\vect X)$ is
small in population mean square, or even that its variance is finite. The
results below are exact for the random variable $q(\vect X)$. They describe
the complete $\ell_A(\vect A)$ exactly only if the surface is quadratic on
the relevant support, or approximately when an additional approximation
argument is justified.

An element $b_i$ has units of log expected output per unit of trait $i$;
$H_{ij}$ has units per product of trait units; and $G_{ij}$ has the matching
product of trait units. Their contractions therefore give variances in log
expected output. Changing the common unit in which reproductive output is
counted adds a constant to $\ell_A$ and does not change either variance.

\paragraph{Biological interpretation.}
The slope describes how conditional expected output changes near the
population mean. Curvature describes how that slope changes. Neither is a
variance until combined with the population distribution. A small curvature
coefficient can contribute appreciable variation if the population occupies
a sufficiently broad range, and a substantial coefficient can contribute
little variation in a narrow population.

\subsection{Why the multivariate Taylor formula has this form}
\label{gaussian:taylor_derivation}

Fix a displacement $\vect x$ whose line segment from $\vect\mu$ to
$\vect\mu+\vect x$ lies in the differentiability neighbourhood. Define a
one-variable function along that segment,
\[
f(t)=\ell_A(\vect\mu+t\vect x),\qquad 0\leq t\leq1.
\]
The chain rule supplies the multivariate information:
\begin{align}
f'(t)&=\sum_i\partial_i\ell_A(\vect\mu+t\vect x)x_i,
\nonumber\\
f''(t)&=\sum_{i,j}\partial_i\partial_j
               \ell_A(\vect\mu+t\vect x)x_i x_j.
\end{align}
Integrate $f''$ once to obtain $f'(s)=f'(0)+\int_0^s f''(t)\,dt$,
then integrate over $s$ from zero to one. Reversing the order of integration
over the triangle $0\leq t\leq s\leq1$ gives
\[
f(1)=f(0)+f'(0)+\int_0^1(1-t)f''(t)\,dt.
\]
Writing $\mat H(\vect a)=\nabla^2\ell_A(\vect a)$ and subtracting the
constant-Hessian part of this integral yields the exact remainder
\begin{equation}
R(\vect x)
 =\int_0^1(1-t)\vect x^{\mathsf T}
       \{\mat H(\vect\mu+t\vect x)-\mat H(\vect\mu)\}
       \vect x\,dt.
\label{gaussian:second_order_remainder}
\end{equation}
The integral of $1-t$ is $1/2$, which explains the coefficient in
Equation~\eqref{gaussian:quadratic_definition}. Continuity of the Hessian
at $\vect\mu$ makes its difference tend uniformly to zero along a
shrinking segment. Multiplication by $\norm{\vect x}^2$ then gives the
stated $o(\norm{\vect x}^2)$ remainder.

If a third derivative also exists continuously, one more integration gives
\begin{equation}
R(\vect x)
 =\frac12\int_0^1(1-t)^2
  D^3\ell_A(\vect\mu+t\vect x)
     [\vect x,\vect x,\vect x]\,dt.
\label{gaussian:integral_remainder}
\end{equation}
Here $D^3\ell_A[\vect x,\vect x,\vect x]$ abbreviates
$\sum_{i,j,k}\partial_i\partial_j\partial_k\ell_A\,x_i x_j x_k$.
If the absolute value of this trilinear form is at most
$M\norm{\vect x}^3$ along the segment, then
\begin{equation}
|R(\vect x)|\leq\frac M2\norm{\vect x}^3
                    \int_0^1(1-t)^2\,dt
                 =\frac M6\norm{\vect x}^3.
\label{gaussian:remainder_bound}
\end{equation}
The derivative bound must hold where the segment runs. A bound established
only in a ball of radius $a_0$ controls
$\E[R(\vect X)^2\,\mathbf1\{\norm{\vect X}\leq a_0\}]$, but leaves
$\E[R(\vect X)^2\,\mathbf1\{\norm{\vect X}>a_0\}]$ uncontrolled.
Because a nondegenerate Gaussian places positive probability outside every
finite ball, a population approximation needs an additional tail argument
or a global derivative bound. We do not replace that requirement with the
observation that a local quadratic fit is close near the mean.

\section{Deriving the Gaussian moments}
\label{gaussian:moments}

Assume first that $\vect X\sim\mathcal N(\vect0,\mat G)$ with
$\mat G\succ0$. Section~\ref{gaussian:singular_support} removes the
positive-definiteness restriction by working only in directions that vary.
We need moments through order four because squaring the quadratic term
produces products of four breeding-value deviations.

The moment-generating function is
\[
M(\vect t)=\E[\exp(\vect t^{\mathsf T}\vect X)].
\]
It is useful because differentiating the exponential with respect to $t_i$
brings down a factor $X_i$. Completing the square in the normal density gives
\begin{align}
-\frac12\vect x^{\mathsf T}\mat G^{-1}\vect x
 +\vect t^{\mathsf T}\vect x
 &=-\frac12(\vect x-\mat G\vect t)^{\mathsf T}
          \mat G^{-1}(\vect x-\mat G\vect t)
   +\frac12\vect t^{\mathsf T}\mat G\vect t.
\end{align}
To check the equality, expand the first quadratic on the right: its two
cross terms each contribute $\vect t^{\mathsf T}\vect x/2$, and its
constant term cancels half of $\vect t^{\mathsf T}\mat G\vect t$.
The integral of the shifted normal density is one. Consequently,
\begin{equation}
M(\vect t)=\exp\!\left(\frac12\vect t^{\mathsf T}\mat G\vect t\right).
\label{gaussian:mgf}
\end{equation}
Differentiation under this Gaussian integral is legitimate: on a bounded
neighbourhood of $\vect t$, a polynomial times
$\exp(\vect t^{\mathsf T}\vect x)$ is dominated by an integrable Gaussian
with a slightly smaller positive precision. In particular, all moments
used here are finite.

For the derivatives, write $u_i=(\mat G\vect t)_i$. Then
$\partial u_i/\partial t_j=G_{ij}$, and
\begin{align}
\partial_i M&=u_iM,\nonumber\\
\partial_j\partial_i M&=(G_{ij}+u_i u_j)M,\nonumber\\
\partial_k\partial_j\partial_i M
 &=(G_{ij}u_k+G_{ik}u_j+G_{jk}u_i+u_i u_j u_k)M.
\label{gaussian:mgf_derivatives}
\end{align}
At $\vect t=\vect0$, every $u_i$ is zero and $M=1$. Therefore
\begin{equation}
\E[X_i]=0,
\qquad \E[X_iX_j]=G_{ij},
\qquad \E[X_iX_jX_k]=0.
\end{equation}
Differentiate the last line of Equation~\eqref{gaussian:mgf_derivatives}
once more. At zero, the derivative of $G_{ij}u_k$ is $G_{ij}G_{kl}$;
the next two terms give $G_{ik}G_{jl}$ and $G_{jk}G_{il}$. Terms retaining
any $u_i$ vanish. We have proved the fourth-moment identity:
\begin{equation}
\E[X_iX_jX_kX_l]
 =G_{ij}G_{kl}+G_{ik}G_{jl}+G_{il}G_{jk}.
\label{gaussian:fourth_moment}
\end{equation}
The three terms correspond to the three ways to pair four positions. This
is the order-four case of the Gaussian product-moment result associated
with \citet{Isserlis1918}; the calculation above establishes the particular
identity needed here rather than treating it as an unexplained rule.
Putting all four indices equal gives $\E[X_i^4]=3G_{ii}^2$.

Subtracting the product of the two second moments yields
\begin{align}
\Cov(X_iX_j,X_kX_l)
 &=\E[X_iX_jX_kX_l]-\E[X_iX_j]\E[X_kX_l]\nonumber\\
 &=G_{ik}G_{jl}+G_{il}G_{jk}.
\label{gaussian:product_covariance}
\end{align}
This covariance of products is the direct input to the quadratic variance.
It includes correlations among traits; independence of the $X_i$ is not
required.

\section{From moments to the variance partition}
\label{gaussian:variance_derivation}

Define the linear term, the uncentred quadratic term, and its centred version
by
\begin{equation}
L=\vect b^{\mathsf T}\vect X,
\qquad Q=\frac12\vect X^{\mathsf T}\mat H\vect X,
\qquad Q_c=Q-\E[Q].
\label{gaussian:terms}
\end{equation}
Since $\E[\vect X]=\vect0$, $\E[L]=0$. For the quadratic mean,
\begin{align}
\E[Q]
 &=\frac12\sum_{i,j}H_{ij}\E[X_iX_j]
   =\frac12\sum_{i,j}H_{ij}G_{ij}\nonumber\\
 &=\frac12\sum_i(\mat H\mat G)_{ii}
   =\frac12\tr(\mat H\mat G).
\label{gaussian:quadratic_mean}
\end{align}
The penultimate equality uses $G_{ij}=G_{ji}$. Thus
\[
q(\vect X)=\E[q(\vect X)]+L+Q_c,
\qquad
\E[q(\vect X)]=\ell_0+\tfrac12\tr(\mat H\mat G).
\]
Centring $Q$ changes its mean but not its variance; it makes the comparison
of the two varying components explicit.

For the linear term,
\begin{align}
\Var(L)
 &=\E\!\left[\left(\sum_i b_iX_i\right)
                   \left(\sum_j b_jX_j\right)\right]\nonumber\\
 &=\sum_{i,j}b_i b_jG_{ij}
   =\vect b^{\mathsf T}\mat G\vect b.
\label{gaussian:linear_variance}
\end{align}
For the quadratic term, covariance is bilinear, so
\begin{align}
\Var(Q_c)
 &=\frac14\sum_{i,j,k,l}H_{ij}H_{kl}
       \Cov(X_iX_j,X_kX_l)\nonumber\\
 &=\frac14\sum_{i,j,k,l}H_{ij}H_{kl}
       (G_{ik}G_{jl}+G_{il}G_{jk}).
\label{gaussian:four_index_variance}
\end{align}
There are two sums in the final line. Interchanging the dummy indices $k$
and $l$ in the first makes it equal to the second because $H_{kl}=H_{lk}$.
For the second sum, regroup the scalar factors:
\begin{align}
\sum_{i,j,k,l}H_{ij}H_{kl}G_{il}G_{jk}
 &=\sum_{i,k}
     \left(\sum_jH_{ij}G_{jk}\right)
     \left(\sum_lH_{kl}G_{li}\right)\nonumber\\
 &=\sum_{i,k}(\mat H\mat G)_{ik}(\mat H\mat G)_{ki}\nonumber\\
 &=\sum_i(\mat H\mat G\mat H\mat G)_{ii}
   =\tr(\mat H\mat G\mat H\mat G).
\end{align}
Again $G_{il}=G_{li}$ was used. No step interchanges the order of two
matrices. The two equal sums and the factor $1/4$ therefore give
\begin{equation}
\Var(Q_c)=\frac12\tr[(\mat H\mat G)^2].
\label{gaussian:quadratic_variance}
\end{equation}
The square is a matrix product, not an elementwise square.

Finally,
\begin{align}
\Cov(L,Q_c)
 &=\E[LQ_c]
   =\frac12\sum_{i,j,k}b_iH_{jk}\E[X_iX_jX_k]\nonumber\\
 &=0.
\label{gaussian:orthogonality}
\end{align}
The constant subtracted in $Q_c$ contributes nothing because $\E[L]=0$.
Equivalently, the linear term is odd under $\vect X\mapsto-\vect X$ and
the quadratic term is even, so their product integrates to zero under a
centrally symmetric Gaussian density.

\begin{proposition}[Gaussian variance partition]
\label{gaussian:variance_partition}
For a centred Gaussian vector with covariance $\mat G\succeq0$, a fixed
vector $\vect b$ and a fixed symmetric matrix $\mat H$, define $q$, $L$
and $Q_c$ by Equations~\eqref{gaussian:quadratic_definition} and
\eqref{gaussian:terms}. Then
\begin{align}
\Vlin&=\Var(L)=\vect b^{\mathsf T}\mat G\vect b,\nonumber\\
\Vquad&=\Var(Q_c)=\frac12\tr[(\mat H\mat G)^2],\nonumber\\
\Var(q)&=\Vlin+\Vquad.
\label{gaussian:variance_formulas}
\end{align}
If the total is positive, the first-order variance share is
\begin{equation}
\Ag=\frac{\Vlin}{\Vlin+\Vquad}.
\label{gaussian:share}
\end{equation}
\end{proposition}
\begin{proof}
For $\mat G\succ0$, the preceding moment calculations establish each
identity; the variance of a sum adds twice the covariance in
Equation~\eqref{gaussian:orthogonality}, which is zero. For singular
$\mat G$, apply the same calculation to its positive-variance coordinates,
as justified in Section~\ref{gaussian:singular_support}.
\end{proof}

Uncorrelated terms need not be independent. For example, if $X$ is standard
normal, $X$ and $X^2-1$ have zero covariance, yet the second is determined
once the first is known. The decomposition above is an orthogonal
decomposition in mean square; it does not describe two independent
biological mechanisms.

\paragraph{Biological interpretation.}
Both components are variation among breeding values in a conditional
expected-fitness function. They are not the additive and epistatic genetic
variance components of the underlying trait, and they are not a partition
of the variance of realised reproductive counts. The environmental and
other variation averaged at each breeding value has already entered the
definition of that conditional mean.

\section{Why the population-fitted line has this share}
\label{gaussian:projection_section}

For a square-integrable response $Y$ and centred predictors $\vect X$, a
population affine fit minimises
\[
J(k,\vect d)=\E[(Y-k-\vect d^{\mathsf T}\vect X)^2].
\]
An affine fit includes both an intercept $k$ and a slope vector $\vect d$.
Expanding this quadratic objective, or differentiating with respect to its
coefficients, gives
\begin{align}
\frac{\partial J}{\partial k}
 &=-2\E[Y-k-\vect d^{\mathsf T}\vect X],\nonumber\\
\nabla_{\vect d}J
 &=-2\E[\vect X(Y-k-\vect d^{\mathsf T}\vect X)].
\end{align}
Setting these derivatives to zero yields the normal equations
\begin{equation}
k=\E[Y],
\qquad
\mat G\vect d=\Cov(\vect X,Y).
\label{gaussian:normal_equations}
\end{equation}
The objective is a convex quadratic in the coefficients. When
$\mat G\succ0$, these equations give its unique minimum. More directly,
if $\widehat Y=k+\vect d^{\mathsf T}\vect X$ satisfies them, its residual
is orthogonal to constants and to every component of $\vect X$. For any
other affine fit $\widehat Y+h$,
\begin{align}
\E[(Y-\widehat Y-h)^2]
 &=\E[(Y-\widehat Y)^2]+\E[h^2]
       -2\E[(Y-\widehat Y)h]\nonumber\\
 &=\E[(Y-\widehat Y)^2]+\E[h^2].
\end{align}
This verifies optimality without relying on a derivative test and shows why
the normal equations are an orthogonal-projection condition.

Apply the equations to $Y=q(\vect X)$. The mean was calculated above, and
\begin{align}
\Cov(\vect X,q)
 &=\Cov(\vect X,\vect b^{\mathsf T}\vect X)
   +\Cov(\vect X,Q_c)\nonumber\\
 &=\mat G\vect b+\vect0.
\end{align}
The last vector is zero component by component by the same third-moment
calculation used in Equation~\eqref{gaussian:orthogonality}. We obtain
\begin{proposition}[Population affine projection]
\label{gaussian:affine_projection}
Under the assumptions of Proposition~\ref{gaussian:variance_partition}, the
best population affine prediction of $q(\vect X)$ is
\begin{equation}
\widehat q_1(\vect X)
 =\ell_0+\frac12\tr(\mat H\mat G)+\vect b^{\mathsf T}\vect X.
\label{gaussian:fitted_affine}
\end{equation}
Its residual is $Q_c$. If $\Var(q)>0$, its population coefficient of
determination is
\begin{equation}
R_1^2
 =1-\frac{\E[(q-\widehat q_1)^2]}{\Var(q)}
 =\frac{\Vlin}{\Vlin+\Vquad}
 =\Ag.
\label{gaussian:regression_share}
\end{equation}
\end{proposition}
\begin{proof}
The normal equations give the displayed intercept and $\vect d=\vect b$
on the varying support. Subtracting the fit from $q$ leaves $Q_c$.
The residual has zero mean and variance $\Vquad$, whereas
$\Var(q)=\Vlin+\Vquad$. Substitution gives the stated $R^2$.
\end{proof}

The tangent plane at the population mean is
$\ell_0+\vect b^{\mathsf T}\vect X$. Its slope agrees with the population
fit under these assumptions, but its intercept omits the mean curvature
contribution. For one trait with $H<0$, the fitted intercept is lower by
$|H|G/2$. Comparing the curve to its tangent without this intercept
adjustment does not give the fitted regression residual.

If $\Vlin>0$, the correlation expression follows separately:
\begin{align}
\Corr(L,L+Q_c)^2
 &=\frac{\Cov(L,L+Q_c)^2}
         {\Var(L)\Var(L+Q_c)}\nonumber\\
 &=\frac{\Vlin^2}{\Vlin(\Vlin+\Vquad)}=\Ag.
\label{gaussian:correlation_identity}
\end{align}
This form requires a nonconstant $L$. The regression form still assigns
$R_1^2=0$ to a constant fit when the target varies.

\paragraph{Biological interpretation.}
Every individual has equal weight in the population mean squared error.
Common breeding values collectively receive more weight because more
individuals have them. A large share therefore means that a linear
description accounts for much of the quadratic target's variation under
those frequencies. It neither identifies gene action nor establishes an
accurate prediction after reproductive reweighting.

\section{What happens when the response includes noise?}
\label{gaussian:noisy_response}

Suppose the observed response has the model
\begin{equation}
Y=q(\vect X)+\kappa+\varepsilon,
\qquad \E[\varepsilon\given\vect X]=0,
\qquad \E[\varepsilon^2]=\sigma_\varepsilon^2<\infty,
\label{gaussian:noise_model}
\end{equation}
where $\kappa$ is constant. Conditional noise variance may depend on
$\vect X$; constant variance and independence are stronger assumptions
than this argument needs. If $f(\vect X)$ is any square-integrable
function for which the product is integrable, conditioning gives
\begin{equation}
\E[f(\vect X)\varepsilon]
 =\E\{f(\vect X)\E[\varepsilon\given\vect X]\}=0.
\label{gaussian:conditional_orthogonality}
\end{equation}
In particular, the residual is orthogonal to the intercept, all linear
predictors, and the complete quadratic conditional mean.

The best quadratic prediction is $q+\kappa$, with residual $\varepsilon$.
The best affine prediction is $\widehat q_1+\kappa$, with residual
$Q_c+\varepsilon$. Hence
\begin{align}
\Var(Y)&=\Vlin+\Vquad+\sigma_\varepsilon^2,\nonumber\\
\E[(Y-\widehat q_1-\kappa)^2]
 &=\Vquad+\sigma_\varepsilon^2.
\end{align}
We can now derive both $R^2$ values rather than assuming they have the same
interpretation as the noiseless share.
\begin{proposition}[The nested-regression ratio]
\label{gaussian:nested_r2}
Under Gaussian predictors, the quadratic target, and
Equation~\eqref{gaussian:noise_model}, let $R_1^2$ and $R_{1+2}^2$ denote
the population coefficients of determination of the affine and complete
quadratic fits. If $\Vlin+\Vquad>0$, then
\begin{align}
R_1^2&=\frac{\Vlin}
 {\Vlin+\Vquad+\sigma_\varepsilon^2},\nonumber\\
R_{1+2}^2&=\frac{\Vlin+\Vquad}
 {\Vlin+\Vquad+\sigma_\varepsilon^2},\nonumber\\
\frac{R_1^2}{R_{1+2}^2}&=\Ag.
\label{gaussian:nested_ratio}
\end{align}
\end{proposition}
\begin{proof}
Subtract the appropriate mean squared residual from $\Var(Y)$ to obtain
each explained variance. Both use the same total response variance.
Dividing cancels that denominator, and the positive quadratic-target
variance ensures that $R_{1+2}^2$ is not zero.
\end{proof}

For $b=2$, $H=-0.5$, $G=1$ and $\sigma_\varepsilon^2=4$, the components
are $4$ and $0.125$. Thus
\[
R_1^2=\frac{4}{8.125}=\frac{32}{65},\qquad
R_{1+2}^2=\frac{4.125}{8.125}=\frac{33}{65},\qquad
\frac{R_1^2}{R_{1+2}^2}=\frac{32}{33}.
\]
Neither individual $R^2$ equals $32/33$. If the quadratic target is
constant but the noise varies, both $R^2$ values are zero and their ratio
is undefined. An infinite residual variance also prevents this variance
calculation.

Chapter~\ref{conditioning:chapter}, especially
Proposition~\ref{conditioning:log_gap_model}, specifies when a logged
individual response can have the conditional mean in
Equation~\eqref{gaussian:noise_model}. The required gap between
$\E[\log W\given\vect A]$ and $\log\E[W\given\vect A]$ must be
constant in breeding value, and the log error must have a finite second
moment. Those are additional observation-model conditions. A general
fitness response, particularly a count with zeros, does not satisfy them
merely because its expected value is positive.

\paragraph{Biological interpretation.}
A ratio of two fit statistics can recover the share only if both fits
target the same conditional relationship with the stated residual model.
The algebra does not turn variation in realised reproduction into
breeding-value variance, and it does not justify logging individual counts.
Directly modelling expected output remains a separate route.

\section{Singular covariance and the boundary cases}
\label{gaussian:singular_support}

A covariance matrix can be singular when some recorded trait combinations
do not vary. An inverse in all recorded coordinates would then be
meaningless, but the variance calculation remains well defined.
Let $r=\operatorname{rank}(\mat G)$ and choose orthonormal columns of an
$n\times r$ matrix $\mat U$ spanning its positive-eigenvalue subspace.
Then
\begin{equation}
\mat G=\mat U\mat\Lambda\mat U^{\mathsf T},
\qquad \mat U^{\mathsf T}\mat U=\mat I_r,
\qquad \mat\Lambda\succ0.
\label{gaussian:support_decomposition}
\end{equation}
If $\vect v$ lies in the null space of $\mat G$, then
$\Var(\vect v^{\mathsf T}\vect X)=0$ and its mean is zero. Therefore
$\vect v^{\mathsf T}\vect X=0$ almost surely. It follows that
$\vect X=\mat U\vect Z$ almost surely, where
$\vect Z=\mat U^{\mathsf T}\vect X\sim\mathcal N(\vect0,\mat\Lambda)$.
In these occupied coordinates,
\begin{equation}
\vect b_r=\mat U^{\mathsf T}\vect b,
\qquad \mat H_r=\mat U^{\mathsf T}\mat H\mat U,
\qquad
q=\ell_0+\vect b_r^{\mathsf T}\vect Z
       +\tfrac12\vect Z^{\mathsf T}\mat H_r\vect Z.
\end{equation}
All positive-definite derivations apply to $\mat\Lambda$. Substitution and
cyclicity of trace recover the formulas in the original coordinates.
The normal equation only requires
$\mat G(\vect d-\vect b)=\vect0$: slope coefficients outside the occupied
subspace are not uniquely identified, but fitted values are unique because
those extra directions contribute zero almost surely.

The boundary distinctions are consequently precise:
\begin{itemize}
\item If $\Vlin=0<\Vquad$, the fitted response is constant and its $R^2$
is zero. The correlation in Equation~\eqref{gaussian:correlation_identity}
is undefined because its first argument has zero variance.
\item If $\Vquad=0<\Vlin$, the quadratic target is affine on the occupied
support and $\Ag=1$.
\item If both components are zero, the target is constant almost surely;
the share and target-regression $R^2$ are undefined, not zero.
\end{itemize}
In particular, $\Vlin=0$ means that the gradient has no component in a
varying direction. It need not mean that every component of the gradient
in an enlarged, partly unoccupied trait space is zero.

\section{What changes outside a Gaussian population?}
\label{gaussian:non_gaussian}

The conditioning and reproductive-weighting definitions do not require
Gaussian breeding values. The moment formulas do. It is useful to derive
the corrections explicitly so that a departure from normality has a
mathematical meaning rather than serving as a general caveat.

Let $\vect X$ now be any centred vector with finite fourth moments. Retain
$\mat G=\E[\vect X\vect X^{\mathsf T}]$, and define the third central
moment tensor and fourth cumulant tensor by
\begin{align}
M_{ijk}&=\E[X_iX_jX_k],\nonumber\\
K_{ijkl}&=\E[X_iX_jX_kX_l]
 -G_{ij}G_{kl}-G_{ik}G_{jl}-G_{il}G_{jk}.
\label{gaussian:cumulants}
\end{align}
The tensor $K$ records the part of the fourth moment that the Gaussian
pairings do not explain. The linear variance and quadratic mean still use
only second moments, so their previous formulas remain valid. Repeating
the covariance calculation, however, gives
\begin{align}
\Cov(L,Q_c)
 &=\frac12\sum_{i,j,k}b_iH_{jk}M_{ijk},\nonumber\\
\Var(Q_c)
 &=\frac12\tr[(\mat H\mat G)^2]
   +\frac14\sum_{i,j,k,l}H_{ij}H_{kl}K_{ijkl}.
\label{gaussian:non_gaussian_corrections}
\end{align}
These are exact identities under the stated finite-moment assumption.
The total variance includes twice the first line in addition to the
linear and actual quadratic variances. Central symmetry makes $M_{ijk}=0$
but need not make $K_{ijkl}=0$.

The slope of the best affine fit can also change. Set
\begin{equation}
u_i=\Cov(X_i,Q_c)=\frac12\sum_{j,k}H_{jk}M_{ijk}.
\end{equation}
On a positive-definite occupied support,
\begin{align}
\Cov(\vect X,q)&=\mat G\vect b+\vect u,\nonumber\\
\vect d&=\vect b+\mat G^{-1}\vect u,\nonumber\\
\Var(\vect d^{\mathsf T}\vect X)
 &=\vect b^{\mathsf T}\mat G\vect b
   +2\vect b^{\mathsf T}\vect u
   +\vect u^{\mathsf T}\mat G^{-1}\vect u.
\label{gaussian:non_gaussian_projection}
\end{align}
Thus the variance of the local linear term and the variance explained by
the best affine fit are different quantities in general. One can still
calculate the actual population regression using the observed
distribution; one cannot infer that its $R^2$ equals the Gaussian
expression simply by substituting an observed covariance.

\begin{example}[Symmetry preserves orthogonality but not the fourth moment]
Let $X=-1$ or $1$ with equal probabilities, $b=1$ and $H=-1$. Then
$G=1$, $\E[X^3]=0$ and $\E[X^4]=1$, so the fourth cumulant is
$1-3=-2$. Since $X^2=1$ always, $Q_c=-(X^2-1)/2=0$. The quadratic
variance is exactly zero. Equation~\eqref{gaussian:non_gaussian_corrections}
gives the same answer, $1/2+(-2)/4=0$. The response is perfectly affine
on these two points, with $R^2=1$, whereas substituting $G=1$ into the
Gaussian formula would give $1/(1+1/2)=2/3$.
\end{example}

\begin{example}[Asymmetry changes the fitted slope]
Let $X=-1$ with probability $2/3$ and $X=2$ with probability $1/3$.
Then $\E[X]=0$, $G=2$, $\E[X^3]=2$, and $\E[X^4]=6$.
Choose $b=1$ and $H=-1/2$, so $q=X-X^2/4$. The centred quadratic term
is $Q_c=-(X^2-2)/4$. Direct averaging gives
\[
\Var(L)=2,\qquad \Var(Q_c)=\frac18,\qquad
\Cov(L,Q_c)=-\frac12.
\]
Consequently $\Var(q)=2+1/8-1=9/8$. Here
$u=H\E[X^3]/2=-1/2$, and the fitted slope is
$d=1+(-1/2)/2=3/4$. The intercept is $\E[q]=-1/2$.
The fitted line $-1/2+3X/4$ equals $q$ at both support points, so its
$R^2$ is one. The Gaussian expression instead uses quadratic variance
$H^2G^2/2=1/2$ and would report $2/(2+1/2)=0.8$.
\end{example}

\paragraph{Biological interpretation.}
Changes in breeding-value frequencies can affect third and fourth moments
as well as covariance. The Gaussian formulas isolate a tractable case,
not a property of all populations with the same covariance. A finite
multilocus simulation can assess the approximation in its specified
design; it cannot prove Gaussianity or make higher moments irrelevant.

\section{Connection to executable checks}

The groups \codepath{gaussian_projection},
\codepath{noisy_regression_ratio}, and
\codepath{boundaries_and_non_gaussian_cases} in
\codepath{code/reproduce_result_checks.py} independently
integrate or fit small specified populations. They check the displayed
values and implementations. The general results above follow from their
proofs, not from passing a finite set of examples. Sample correlations
from a finite Gaussian sample can differ from the population identity
because the sample's third and fourth moments are not exactly their
population values.
